\chapter{Hamiltonian, Boltzmann, to Navier stokes}

Sources:
\begin{itemize}
  \item https://www.unm.edu/~aierides/505/bbgky3.pdf?
  \item https://journals.aps.org/rmp/pdf/10.1103/RevModPhys.52.569
\end{itemize}

\begin{itemize}
  \item Introduction
  \item Scope
  \item Brief overview of the three
  \item Liouville to BBGKY
  \item The BBGKY Hierarchy
  \item The Boltzmann equation
  \item The H-Theorem
  \item Boltzmann to Navier-Stokes
  \item Hydro limit
\end{itemize}

\section{Introduction}

Hello, and welcome to another video. On the connection between Hamiltonian dynamics, Boltzmann equation, and navier-stokes

\section{Scope}

Here I plan to derive the Navier-stokes using the Boltzmann, and the Boltzmann using the N-body Hamiltonian. Then using how we did that, examine their connection and when it is appropriate to use one and not the other.


\section{Brief Overview}
....


\section{From Liouville to BBGKY}

Taking an N-body Hamiltonian

\[
H = \frac{1}{2m}\sum^N_{i=1} \va{p_i}^2 + \sum^N_{i=1} V\va{r_i} + \sum^N_{i<j} U(\va{r_i} - \va{r_j})
  \]

  Where the Hamiltonian constains an external force $\va{F} = - \nabla V$, that acts equally on every particle. There is also two-body interaction between particles captured by the potential energy term. Later on we will make the assumption that this potential is short ranged. I will also be notating position as r rather than q for convenience.

For our interests, we wish to examine it as a probobility distribution, like all probobilitity distributions they are normalized as follows
\[
\int dV \, f(\va{r},\va{p},t) = 1 \quad dV = \prod_{i=1}^N d^3 r_i d^3 p_i
  \]

  Further, since the probability is locally conserved, the equation most follow continuity. Any probability change in one point in phase-space must by its nature be compensated by neighboring points.

  This fairly trivially corresponds to the equation
\[
\pdv{f}{t} +\pdv{}{\va{r_i}} \cdot (\dot{\va{r_i}}f) +\pdv{}{\va{p_i}} \cdot (\dot{\va{p_i}}f) = 0
  \]

  Using the nature of what a Hamiltonian and its derivatives we can transform this equation

  \[
\pdv{f}{t} +\pdv{}{\va{r_i}} \cdot (\pdv{H}{\va{p}} f) +\pdv{}{\va{p_i}} \cdot (\pdv{H}{\va{r_i}}f) = 0 \longrightarrow
  \]
\[
\pdv{f}{t} + \sum_{i=1}^N \pdv{f}{\va{r_i}} \cdot \pdv{H}{\va{p_i}} + \pdv{f}{\va{p_i}} \cdot \pdv{H}{\va{r_i}} = 0
  \]

This final equation is the Liouville’s equation. It is the statement that probability
doesn’t change as you follow it along any trajectory in phase space, as is seen by
writing the Liouville equation as a total derivative

\[
\dv{f}{t} = \pdv{f}{t} + \sum_{i=1}^N \pdv{f}{\va{r_i}} \cdot \pdv{H}{\va{p_i}} + \pdv{f}{\va{p_i}} \cdot \pdv{H}{\va{r_i}} = 0
  \]

  we can better understand this equation using the familiar Liouville theorem that we have used in the past to express in poisson brackets defined at
\[
\pb{A}{B} \equiv \pdv{A}{\va{r_i}} \cdot \pdv{B}{\va{p_i}} + \pdv{A}{\va{p_i}} \cdot \pdv{B}{\va{r_i}}
  \]

  With this definition one can clearly see the above equation finds that

\[
\pdv{f}{t} = \pb{H}{f}
  \]

  Meaning that the equilibrium distribution is explicitly time independent.

  To make this statement true, we must make $f$ a function of $H$, the most famous example being Boltzmann's distribution $f \sim e^{-\beta H}$, there is nothing that requires this as there are several such solutions given this limited definition.

Now, let us add a new function $A(\va{r}_i,\va{p}_i)$ on the phase space.With the expectation value

\[
\langle A \rangle = \int dV A(\va{r}_i,\va{p}_i) f(\va{r}_i,\va{p}_i, t)
  \]

  We can then find the time dependence of the expectation value of A by

\begin{align}
  \dv{\langle A \rangle}{t} \, &= \, \int \, dV \, A \pdv{f}{t} \\
                               &= \, \int \, dV \, A \left(\pdv{f}{\va{r_i}} \cdot \pdv{H}{\va{p_i}} + \pdv{f}{\va{p_i}} \cdot \pdv{H}{\va{r_i}} \right) \\
                              &= \, \int \, dV \, \left(\pdv{A}{\va{r_i}} \cdot \pdv{H}{\va{p_i}} + \pdv{A}{\va{p_i}} \cdot \pdv{H}{\va{r_i}} \right) f
\end{align}

where we hav0e integrated by parts to get to the last line, throwing away boundary terms which is justified in this context because f is normalized which ensures that we must have $f \to 0$ in asymptotic parts of phase space. Finally, we learn that
\[
\dv{\langle A \rangle}{t} = \int dV \, \pb{A}{H}f = \langle \pb{A}{H} \rangle
\]

This looks very closely to the quantum expectation values.

\subsection{The BBGKY Hierarchy}

Now, that we have it in a distribution formulism, we still have a problem, we still have a function with an obscene amount of variables, assuming we have N $\approx 10^{23}$, so we will focus on one particles, but since all of the particles are identical, we aren't actually privileging one particle over the others. From we find that a series of equations; energy flux, average density, average velocity, etc all require a function $f_1$, we define $f_1$ time derivative as:

\[
\pdv{f_1}{t} = N \int \prod_{i=1}^n d^3r_i d^3 p_i \pb{H}{f}
  \]

  Using the Hamiltonian from earlier this becomes
\[
  \frac{\partial f_1}{\partial t} = N \int \prod_{i=2}^{N} d^3r_i\, d^3p_i \left( -\frac{\mathbf p}{m}\cdot\frac{\partial f}{\partial \mathbf r} +\frac{\partial V(\mathbf r)}{\partial \mathbf r}\cdot\frac{\partial f}{\partial \mathbf p} +\sum_{k=2}^{N} \frac{\partial U(\mathbf r-\mathbf r_k)}{\partial \mathbf r}\cdot \frac{\partial f}{\partial \mathbf p} \right)
\]

Now, whenever $j =2,...N$, we can always integrate by parts to move the derivatives
away from f and onto the other terms. And, in each case, the result is simply zero
because when the derivative is with respect to $\va{r}_j$ , the other terms depend only on $\va{p}_i$
and vice-versa. We’re left only with the terms that involve derivatives with respect
to $\va{r}_1$ and $\va{p}_1$ because we can’t integrate these by parts. Let’s revert to our previous
notation and call $\va{r}_1 \equiv \va{r}$ and $\va{p}_1 \equiv \va{p}$. We have

\begin{align}
  \frac{\partial f_1}{\partial t} &= N \int \prod_{i=2}^{N} d^3r_i\, d^3p_i \left( -\frac{\mathbf p}{m}\cdot\frac{\partial f}{\partial \mathbf r} +\frac{\partial V(\mathbf r)}{\partial \mathbf r}\cdot\frac{\partial f}{\partial \mathbf p} +\sum_{k=2}^{N} \frac{\partial U(\mathbf r-\mathbf r_k)}{\partial \mathbf r}\cdot \frac{\partial f}{\partial \mathbf p} \right) \\
 &= \{H_1,f_1\} + N \int \prod_{i=2}^{N} d^3r_i\, d^3p_i \sum_{k=2}^{N} \frac{\partial U(\mathbf r-\mathbf r_k)}{\partial \mathbf r}\cdot \frac{\partial f}{\partial \mathbf p}. \tag{2.10}
\end{align}


We find this expression can be defined by a Liouville-like equation with an added collision term.
\[
\pdv{f_1}{t} = \pb{H_1}{f_1} + \left(\pdv{f_1}{H}  \right)_{coll}
  \]

  The first term is sometimes referred to as the streaming term. It tells you how the
particles move in the absence of collisions. The second term, known as the collision
integral, is given by the second term in In fact, because all particles are the
same, each of the $(N-1)$ terms in $\sum^N_{k=2}$ from earlier are identical and we can write

.... continue


\section{Motivating the Boltzmann Equation}

Now that we have the structure of the Boltzmann equation
\[
\pdv{f_1}{t} = \pb{H_1}{f_1} + \left(\pv{f_1}{t}  \right)
  \]

  but now we need to create an expression for the collision integral in terms of $f_1$.
